\chapter{Particle-Flow Particle Filters and Proposal Correction}
\label{ch:bf-dpf-pfpf}

The previous chapter developed particle flow as a transport construction.  The
main unresolved question left there was whether the transported cloud itself may
already be treated as the corrected filtering object.  In general, the answer is
no.  Outside exact special cases, the flow should be interpreted first as a
proposal mechanism.  The role of the present chapter is to make that proposal
interpretation explicit and to derive the corresponding importance correction.

This chapter is therefore the mathematical bridge between raw flow transport and
a more principled particle-based likelihood interpretation.  It does not yet
introduce differentiable resampling or optimal transport relaxations.  Instead,
it asks a prior question: once a particle has been moved by a deterministic flow
map, what density does that transported particle actually have, and what
importance weight is required to target the intended posterior law?

\section{Objects, assumptions, and source roles}
\label{sec:bf-pfpf-objects}

The PF-PF argument depends on naming the target and proposal objects before
writing any weight.  This chapter uses:

\begin{center}
\small
\begin{tabular}{p{0.24\textwidth} p{0.66\textwidth}}
\toprule
Object & Meaning \\
\midrule
$x_{t-1}$ & ancestor state after the previous filtering step \\
$x_{t,0}$ & pre-flow particle sampled from a one-step proposal \\
$x_{t,1}$ & post-flow particle after applying the flow map \\
$q_{t,0}$ & pre-flow proposal density for $x_{t,0}$ conditional on
$(x_{t-1},y_t)$ \\
$\Phi_t$ & differentiable flow map from pre-flow to post-flow coordinates \\
$\Phi_t^{-1}$ & inverse map used to express the pre-image of a post-flow
particle \\
$J_t=D\Phi_t$ & forward Jacobian matrix of the flow map \\
$q_{t,1}$ & induced post-flow proposal density for $x_{t,1}$ \\
$\gamma_t$ & unnormalized one-step filtering target density \\
$\tilde w_t$ & unnormalized importance weight correcting $q_{t,1}$ to
$\gamma_t$ \\
\bottomrule
\end{tabular}
\end{center}

The source role is narrow.  General SMC references
\citep{doucet2001sequential,chopin2020introduction} supply the proposal and
importance-weight framework.  Invertible particle-flow particle filtering
\citep{li2017particle} motivates treating particle flow as a proposal
transformation with a Jacobian correction.  These sources do not remove the
need to state the target/proposal ratio locally, and they do not imply that
finite-particle, numerical, or HMC target issues have disappeared.

\section{Why proposal correction is needed}
\label{sec:bf-pfpf-why-correction}

In the linear-Gaussian special case studied in the previous chapter, EDH can
recover the exact Kalman update.  Outside that case, however, the transport map
is typically built under Gaussian closure or local linearization.  Even if the
transport moves particles toward high posterior density, the moved cloud is not
automatically distributed according to the exact filtering law.

The right interpretation is therefore proposal-theoretic.  Let
$x_{t,0}$ denote a pre-flow particle sampled from a proposal law, and let the
flow map transport it to
\begin{equation}
\label{eq:bf-pfpf-flow-map}
  x_{t,1} = \Phi_t(x_{t,0}).
\end{equation}
The question is then not merely where the particle moved, but which density on
$x_{t,1}$ is induced by the transformation and how that induced proposal density
must be corrected to target the filtering posterior.

This is the basic PF-PF point emphasized in the invertible particle-flow
literature: the flow is best viewed as a proposal-improving device until its
proposal-density effect has been accounted for.  The exact change-of-variables
identity is not itself the approximation; the approximation enters through the
choice of flow family, the numerical integration of the flow, and the finite
particle system.

\section{Change of variables under the flow map}
\label{sec:bf-pfpf-change-of-variables}

Assume that for each fixed time step and observation, the flow map
$\Phi_t: \mathbb{R}^{n_x} \to \mathbb{R}^{n_x}$ is a differentiable bijection on
the region of interest and that its forward Jacobian
$D\Phi_t(x_{t,0})$ is an $n_x\times n_x$ nonsingular matrix along the proposal
paths being weighted.  Equivalently, the calculation below is a local
diffeomorphism calculation on the support actually used by the proposal.  Let
the pre-flow proposal density be
\begin{equation}
\label{eq:bf-pfpf-preflo-density}
  q_{t,0}(x_{t,0} \mid x_{t-1}, y_t).
\end{equation}
Then the transported particle density is given by the standard
change-of-variables formula,
\begin{equation}
\label{eq:bf-pfpf-postflow-density}
  q_{t,1}(x_{t,1} \mid x_{t-1}, y_t)
  = q_{t,0}(x_{t,0} \mid x_{t-1}, y_t)
    \left|\det D\Phi_t(x_{t,0})\right|^{-1},
\end{equation}
where $x_{t,0} = \Phi_t^{-1}(x_{t,1})$.
The determinant in \eqref{eq:bf-pfpf-postflow-density} is the determinant of
the \emph{forward} map from $x_{t,0}$ to $x_{t,1}$.  Equivalently,
\[
  q_{t,1}(x_{t,1}\mid x_{t-1},y_t)
  =
  q_{t,0}(\Phi_t^{-1}(x_{t,1})\mid x_{t-1},y_t)
  \left|\det D\Phi_t^{-1}(x_{t,1})\right|.
\]
These two forms are identical because
$|\det D\Phi_t^{-1}(x_{t,1})|
=|\det D\Phi_t(x_{t,0})|^{-1}$.  The chapter uses the forward-map convention
because the flow ODE naturally evolves $x_{t,0}$ to $x_{t,1}$ and the
log-determinant ODE below accumulates $\log|\det D\Phi_t(x_{t,0})|$.

Thus the Jacobian determinant is not an optional numerical decoration.  It is
part of the proposal density itself.  Any flow-based importance scheme that
forgets this factor changes the proposal law and therefore changes the targeting
relationship of the resulting importance weights.

\section{Proposal-corrected PF-PF weights}
\label{sec:bf-pfpf-corrected-weights}

Let the intended one-step filtering target be proportional to
\begin{equation}
\label{eq:bf-pfpf-target-density}
  \gamma_t(x_t)
  \propto p_\theta(y_t \mid x_t)
          p_\theta(x_t \mid x_{t-1}).
\end{equation}
This target is conditional on the ancestor state.  In a full particle filter,
the ancestor itself is random and selected from the previous weighted cloud.
The present derivation is therefore a one-step conditional density-ratio
calculation; the full likelihood estimator also depends on ancestor selection,
resampling policy, particle count, and time recursion.

For a generic pre-flow proposal $q_{t,0}$, the corrected importance weight is
\begin{equation}
\label{eq:bf-pfpf-generic-weight}
  \tilde w_t
  \propto
  \frac{\gamma_t(x_{t,1})}
       {q_{t,1}(x_{t,1}\mid x_{t-1},y_t)}
  =
  \frac{p_\theta(y_t\mid x_{t,1})
        p_\theta(x_{t,1}\mid x_{t-1})
        \left|\det D\Phi_t(x_{t,0})\right|}
       {q_{t,0}(x_{t,0}\mid x_{t-1},y_t)}.
\end{equation}
The bootstrap/prior-proposal formula below is the special case
$q_{t,0}=p_\theta(\cdot\mid x_{t-1})$.

If the pre-flow particle is sampled from the transition prior,
\begin{equation}
\label{eq:bf-pfpf-preflo-prior}
  q_{t,0}(x_{t,0} \mid x_{t-1}, y_t)
  = p_\theta(x_{t,0} \mid x_{t-1}),
\end{equation}
and then transported to $x_{t,1}=\Phi_t(x_{t,0})$, the post-flow proposal
density is
\begin{equation}
\label{eq:bf-pfpf-postflow-prior}
  q_{t,1}(x_{t,1} \mid x_{t-1}, y_t)
  = p_\theta(x_{t,0} \mid x_{t-1})
    \left|\det D\Phi_t(x_{t,0})\right|^{-1}.
\end{equation}
The corresponding unnormalized importance weight is therefore
\begin{equation}
\label{eq:bf-pfpf-weight}
  \tilde w_t
  \propto
  \frac{p_\theta(y_t \mid x_{t,1})
        p_\theta(x_{t,1} \mid x_{t-1})}
       {q_{t,1}(x_{t,1} \mid x_{t-1}, y_t)}
  =
  \frac{p_\theta(y_t \mid x_{t,1})
        p_\theta(x_{t,1} \mid x_{t-1})
        \left|\det D\Phi_t(x_{t,0})\right|}
       {p_\theta(x_{t,0} \mid x_{t-1})}.
\end{equation}
Equation \eqref{eq:bf-pfpf-weight} is the central correction formula for the
PF-PF viewpoint.  It shows exactly what proposal correction restores: once the
transformation of proposal density has been accounted for, the flow-transported
particle may be reweighted against the desired posterior factor in the usual
importance-sampling way.

The formula also shows exactly what is \emph{not} optional in code.  A PF-PF
implementation must evaluate, on the same path, the observation density at the
post-flow particle, the transition density at the post-flow particle, the
pre-flow proposal density at the pre-image, and the forward log determinant.
Changing any one of these objects changes the proposal-to-target ratio being
computed.

\section{EDH/PF and LEDH/PF}
\label{sec:bf-pfpf-edh-ledh}

The correction formula is structurally the same whether the flow map comes from
EDH or LEDH.  The difference lies in how the map itself is built.

\subsection{EDH/PF}

For EDH/PF, the transport map is generated by the global affine ODE from the
previous chapter.  The resulting proposal is easier to analyze because the flow
field is common across the cloud.  The main approximation point is the Gaussian
closure used to derive the affine coefficients.  If the affine ODE is integrated
with a common map $\Phi_t$, the same forward log determinant applies to every
particle up to its path through the time-varying affine field.

\subsection{LEDH/PF}

For LEDH/PF, the map is particle-specific: each particle evolves under a local
linearization and therefore under its own drift coefficients.  This can improve
local adaptation to nonlinear observation structure, but it also makes the
proposal analysis more delicate.  The local map remains proposal-correctable in
principle through the same change-of-variables logic, but the numerical burden
of evaluating particle-specific Jacobians or log-determinants becomes more
significant.  A code path that reuses an EDH log determinant for LEDH has
changed the proposal density unless the local maps are actually identical.

The important mathematical point is that EDH/PF and LEDH/PF are not merely flow
algorithms.  They are proposal-corrected particle methods.  That is why they
occupy a different target-status category from raw flow transport.

\section{Li--Coates Algorithm 1 as an implementation contract}
\label{sec:bf-pfpf-li-coates-algorithm-contract}

Algorithm~1 of \citet{li2017particle} is more specific than the generic
proposal-correction derivation above.  It is a recursive particle filter in
which each particle carries both a state and a covariance object.  For
BayesFilter, this algorithm must be read as an implementation contract, not as
a vague instruction to apply a deterministic affine map.

At the start of time step $k$, particle $i$ has an ancestor
$x_{k-1}^i$, an importance weight $w_{k-1}^i$, and a posterior covariance
$P_{k-1}^i$ inherited from the previous time step.  The covariance lifecycle in
Algorithm~1 is
\begin{equation}
\label{eq:bf-pfpf-alg1-covariance-lifecycle}
  (x_{k-1}^i,P_{k-1}^i)
  \xrightarrow{\mathrm{EKF/UKF\ prediction}}
  (m_{k|k-1}^i,P^i)
  \xrightarrow{\mathrm{EKF/UKF\ update}}
  (m_{k|k}^i,P_k^i).
\end{equation}
The source states that EKF or UKF covariance prediction equations may be used
when the dynamic model is nonlinear.  The source does not make optimal
transport resampling part of Algorithm~1, and it should not be cited as
evidence for any BayesFilter OT relaxation.  In the BayesFilter rebuild, using
UKF for the prediction and update arrows in
\eqref{eq:bf-pfpf-alg1-covariance-lifecycle} is a reviewed implementation
choice requested for this program; it is not a claim that the paper's reported
experiments used UKF in every scenario.

After the prediction arrow, Algorithm~1 defines two distinct state objects:
\begin{align}
\label{eq:bf-pfpf-alg1-zero-noise-anchor}
  \bar\eta_0^i &= g_k(x_{k-1}^i,0),\\
\label{eq:bf-pfpf-alg1-preflo-sample}
  \eta_0^i &= g_k(x_{k-1}^i,v_k).
\end{align}
The first is the zero-noise transition anchor used to drive the local
linearization.  The second is the actual pre-flow proposal particle drawn
through the dynamic model.  The implementation must keep both objects.  If the
code transports only the actual proposal and never maintains the auxiliary
zero-noise path, then the local Jacobian field is not the one specified by
Algorithm~1.

For pseudo-time grid
$0=\lambda_0<\lambda_1<\cdots<\lambda_{N_\lambda}=1$ with
$\epsilon_j=\lambda_j-\lambda_{j-1}$, Algorithm~1 repeats the following
particle-specific update.  Set $P=P^i$ and compute
$A_j^i(\lambda_j)$ and $b_j^i(\lambda_j)$ from the LEDH formulas using the
linearization at the current auxiliary state $\bar\eta^i$.  Then move the
auxiliary and actual states by the same local affine step:
\begin{align}
\label{eq:bf-pfpf-alg1-auxiliary-step}
  \bar\eta^i
  &\leftarrow
  \bar\eta^i
  +\epsilon_j\{A_j^i(\lambda_j)\bar\eta^i+b_j^i(\lambda_j)\},\\
\label{eq:bf-pfpf-alg1-actual-step}
  \eta_1^i
  &\leftarrow
  \eta_1^i
  +\epsilon_j\{A_j^i(\lambda_j)\eta_1^i+b_j^i(\lambda_j)\}.
\end{align}
At the same pseudo-time step, the forward determinant multiplier is accumulated
as
\begin{equation}
\label{eq:bf-pfpf-alg1-theta-product}
  \theta^i
  \leftarrow
  \theta^i
  \left|\det\{I+\epsilon_j A_j^i(\lambda_j)\}\right|,
  \qquad
  \theta^i_0=1.
\end{equation}
Thus, at the end of the pseudo-time integration,
\begin{equation}
\label{eq:bf-pfpf-alg1-det-product}
  \theta^i
  =
  \prod_{j=1}^{N_\lambda}
  \left|\det\{I+\epsilon_j A_j^i(\lambda_j)\}\right|.
\end{equation}
This is the discrete Algorithm~1 version of the forward determinant in
\eqref{eq:bf-pfpf-weight}.  It is particle-specific because the LEDH
coefficients are particle-specific.

After the flow, Algorithm~1 sets $x_k^i=\eta_1^i$ and applies the PF-PF
importance update
\begin{equation}
\label{eq:bf-pfpf-alg1-weight}
  w_k^i
  \propto
  \frac{
    p(x_k^i\mid x_{k-1}^i)
    p(z_k\mid x_k^i)
    \theta^i
  }{
    p(\eta_0^i\mid x_{k-1}^i)
  }
  w_{k-1}^i .
\end{equation}
The denominator is the density of the pre-flow transition sample, while the
transition density in the numerator is evaluated at the post-flow particle
$x_k^i$.  Replacing both densities by the same value, dropping $\theta^i$, or
using a determinant from a different map changes the proposal-to-target ratio.

Finally, the weights are normalized, the EKF/UKF update produces
$P_k^i$, and any resampling step resamples the triple
\begin{equation}
\label{eq:bf-pfpf-alg1-resampling-triple}
  \{x_k^i,P_k^i,w_k^i\}_{i=1}^{N_p}.
\end{equation}
The covariance component is part of the particle state for the next recursion.
Resampling only the states while leaving covariances in their old order breaks
the Algorithm~1 lifecycle.  Likewise, introducing differentiable or optimal
transport resampling may be a BayesFilter extension, but it must be labelled as
an extension and audited separately from the Li--Coates source algorithm.

\subsection{Why the OT block begins here}
\label{sec:bf-pfpf-ot-why-here}

The next chapters do not reopen proposal correction.  By the end of the PF-PF
stage, the object of record is already the weighted post-flow ensemble together
with its carried particle-local state.  The remaining question is different: once
that weighted object exists, what equal-weighting rule should act on it?  The OT
block begins here because BayesFilter still needs a transport object that turns
that weighted PF-PF output into an equally weighted downstream representation.

\subsection{What the later OT chapters inherit from PF-PF}
\label{sec:bf-pfpf-ot-handoff}

The OT survey companion note sharpens the next boundary in a useful way.  After
PF-PF correction, the BayesFilter object handed to a later resampling layer is
not a vague ``particle cloud.''  It is a weighted empirical ensemble together
with any carried auxiliary state.  In the Algorithm~1 reading fixed above, the
state handed forward is precisely the weighted triple
\eqref{eq:bf-pfpf-alg1-resampling-triple}.  Thus a later optimal-transport
layer does not merely need an OT scalar or distance certificate.  It needs a
usable transport object that can act on the weighted ensemble:
\begin{equation}
\label{eq:bf-pfpf-transport-object-inventory}
  \text{coupling matrix, barycentric transform, factorized transport operator,}
  \text{ or directly transported equal-weight cloud.}
\end{equation}
A fast value of \(\operatorname{OT}(\mu,\nu)\) alone is not enough to continue the filtering
recursion, because the filter must actually produce the next equal-weighted
particle representation.

Write the PF-PF output cloud schematically as
\begin{equation}
\label{eq:bf-pfpf-weighted-output-cloud}
  \mu_t^N
  =
  \sum_{i=1}^{N_p} w_k^i\,\delta_{x_k^i},
\end{equation}
with carried covariance or particle-local state attached to each support point.
A deterministic OT equal-weighting layer or entropic OT/Sinkhorn layer, in the
spirit of ensemble-transform and differentiable OT resampling
\citep{reich2013nonparametric,corenflos2021differentiable}, then asks for a
map from the weighted object \(\mu_t^N\) to an equal-weight object.  The next
chapters will spell out that map through couplings, barycentric projections, and
teacher-versus-numerical transport objects.  The present chapter's narrower role
is only to make the handoff precise: PF-PF supplies the weighted source object,
and any later OT resampling route must declare how that weighted object is
actually transported rather than merely reporting a scalar transport score.

These details give a sharper audit checklist than the generic PF-PF equations:
\begin{enumerate}[label=(\roman*)]
  \item each particle has its own $P_{k-1}^i$, predicted $P^i$, and updated
    $P_k^i$;
  \item the local LEDH Jacobian is evaluated along the particle's auxiliary
    zero-noise path;
  \item the actual proposal particle and the auxiliary path are both migrated;
  \item the determinant product in
    \eqref{eq:bf-pfpf-alg1-det-product} uses the same local matrices that
    migrated the actual proposal particle;
  \item the weight uses the post-flow transition density, the observation
    density, the determinant product, the pre-flow proposal density, and the
    previous particle weight;
  \item resampling, if used, resamples covariances together with states and
    weights.
\end{enumerate}

\section{Jacobian determinant and log-determinant evolution}
\label{sec:bf-pfpf-logdet}

Directly forming a full Jacobian matrix for every particle and time step may be
numerically and computationally burdensome.  A more useful identity is obtained
by differentiating the Jacobian matrix along the flow.

Let
\begin{equation}
\label{eq:bf-pfpf-jacobian-matrix}
  J_t(\lambda)
  = \frac{\partial x_t(\lambda)}{\partial x_t(0)}.
\end{equation}
The terminal forward determinant in
\eqref{eq:bf-pfpf-postflow-density} is $\det J_t(1)$.  The inverse-density form
would use $-\log|\det J_t(1)|$ instead.  This sign distinction is one of the
most common PF-PF implementation errors.

If the flow satisfies
\begin{equation}
\label{eq:bf-pfpf-general-flow}
  \frac{dx_t(\lambda)}{d\lambda} = v_{t,\lambda}(x_t(\lambda)),
\end{equation}
then differentiation with respect to the initial point gives
\begin{equation}
\label{eq:bf-pfpf-jacobian-ode}
  \frac{dJ_t(\lambda)}{d\lambda}
  = \nabla v_{t,\lambda}(x_t(\lambda)) J_t(\lambda),
  \qquad
  J_t(0)=I.
\end{equation}
The derivation is direct: differentiate the flow equation with respect to the
initial coordinate $x_t(0)$ and apply the chain rule.  Applying Jacobi's formula
$d\log|\det J|/d\lambda=\operatorname{tr}(J^{-1}dJ/d\lambda)$ then yields
\begin{equation}
\label{eq:bf-pfpf-logdet-ode}
  \frac{d}{d\lambda}
  \log \left|\det J_t(\lambda)\right|
  = \operatorname{tr}
    \bigl(\nabla v_{t,\lambda}(x_t(\lambda))\bigr).
\end{equation}
because $J_t^{-1}(\nabla v_{t,\lambda})J_t$ has the same trace as
$\nabla v_{t,\lambda}$.  Thus
\[
  \log|\det D\Phi_t(x_{t,0})|
  =
  \int_0^1
  \operatorname{tr}
  \bigl(\nabla v_{t,\lambda}(x_t(\lambda))\bigr)\,d\lambda.
\]
For an affine flow
$v_{t,\lambda}(x)=A_t(\lambda)x+b_t(\lambda)$, this simplifies to
\begin{equation}
\label{eq:bf-pfpf-logdet-affine}
  \frac{d}{d\lambda}
  \log \left|\det J_t(\lambda)\right|
  = \operatorname{tr}(A_t(\lambda)).
\end{equation}
The practical consequence is clear: one may integrate a scalar log-determinant
alongside the particle-flow ODE rather than carrying a full dense Jacobian for
every particle in the affine case.  In more general settings, the same identity
still shows that the determinant is part of the dynamical system, not an
external afterthought.  This also gives a direct debugging criterion: if the
corrected weights behave inconsistently with the transported cloud, the first
mathematical suspects are the Jacobian/log-determinant path, the numerical flow
integration, and the finite-particle approximation rather than the formal
change-of-variables identity itself.

\section{What the correction restores, and what it does not}
\label{sec:bf-pfpf-what-restored}

The proposal correction in \eqref{eq:bf-pfpf-weight} restores a clearer
importance-sampling interpretation.  More precisely, it restores the correct
proposal-to-target density ratio for the transported proposal generated by the
flow map.  This is mathematically important: it distinguishes PF-PF from a pure
transport approximation that simply declares the endpoint cloud to be posterior-
like without density correction.

What the correction does \emph{not} automatically restore is an exact analytic
filter in the nonlinear case.  Several approximation layers may still remain:
\begin{itemize}
  \item the flow may itself be derived under Gaussian closure or local
    linearization;
  \item the ODE is integrated numerically rather than in closed form;
  \item the particle system still uses a finite number of particles;
  \item Jacobian or log-determinant quantities may be computed with numerical
    approximation.
\end{itemize}
Therefore PF-PF is mathematically stronger than raw flow transport, but it is
not automatically equivalent to an exact nonlinear filtering likelihood.  The
exact identity in this chapter is the change-of-variables correction for the
proposal density; what remains approximate are the flow construction itself when
closure or local linearization is used, the numerical integration of that flow,
and the finite-particle system used to estimate the target.  Its status is
better described as a proposal-corrected particle approximation whose value-side
interpretation is clearer than that of uncorrected flow transport.

At trajectory level, this one-step correction is multiplied through the
filtering recursion in the same spirit as ordinary sequential importance
sampling.  For a particle path and ancestor history, the product of incremental
corrections defines the path proposal weight.  That statement is still
conditional on using the stated proposal, the stated flow map, and the stated
determinant convention at each time step.  It is not a separate theorem that the
finite-particle likelihood estimate is exact in the nonlinear case.

\section{Implementation audit obligations}
\label{sec:bf-pfpf-audit-obligations}

The formulas above imply concrete checks before PF-PF can be used as a
value-side candidate in later HMC experiments.

\begin{enumerate}[label=(\roman*)]
  \item \textbf{Affine density check}: choose a low-dimensional affine map with
    known determinant and compare \eqref{eq:bf-pfpf-postflow-density} against a
    direct transformed-density calculation.
  \item \textbf{Forward/inverse sign check}: verify that the proposal density
    uses $-\log|\det D\Phi_t|$ while the numerator form in
    \eqref{eq:bf-pfpf-weight} contributes $+\log|\det D\Phi_t|$.
  \item \textbf{Jacobian ODE check}: compare the integrated trace in
    \eqref{eq:bf-pfpf-logdet-ode} with an automatic-differentiation or
    finite-difference determinant in small dimension.
  \item \textbf{Corrected versus uncorrected weights}: run the same transported
    cloud with and without the determinant/proposal-density term to confirm that
    the diagnostic detects the expected discrepancy.
  \item \textbf{Path consistency}: on fixed random seeds, verify that the scalar
    value and any autodiff gradient use the same proposal density, flow map, and
    log-determinant path.
\end{enumerate}

These are not tuning diagnostics.  They are algebraic and numerical checks of
the target/proposal object itself.

\section{Why this chapter matters for the later HMC discussion}
\label{sec:bf-pfpf-hmc-bridge}

This chapter is the first point in the rebuilt DPF block where a sharper target
question becomes possible.  Raw flow transport gives a geometrically appealing
approximation, but PF-PF introduces an explicit proposal density and an explicit
correction factor.  That is why later HMC-suitability analysis will treat
proposal-corrected EDH/PF or LEDH/PF as the first rungs in this ladder with an
explicit proposal-corrected value-side interpretation.  That is a narrower
statement than saying they are HMC-ready.

The key lesson is not that PF-PF solves all target issues.  The key lesson is
that it restores the proposal-correction logic that a valid particle method
requires.  This is exactly the right place in the monograph to separate
flow-as-transport from flow-as-proposal.

\section{What this chapter does and does not claim}
\label{sec:bf-pfpf-chapter-boundary}

This chapter derives the proposal-correction structure for flow-based particle
methods.  It does \emph{not} yet claim:
\begin{itemize}
  \item that differentiable resampling has been resolved;
  \item that the resulting method is already the final HMC target for nonlinear
    DSGE or MacroFinance models;
  \item that OT or learned transport relaxations have been introduced;
  \item that finite-particle, discretization, and Jacobian numerical issues are
    negligible.
\end{itemize}
Those belong to the later resampling and HMC-assessment chapters.  The role of
this chapter is more specific: it explains mathematically how a particle flow is
turned into a proposal-corrected particle method and what that correction does
and does not restore.
